\title{Direct Preference Optimization: Your Language Model is Secretly a Reward Model}

\begin{abstract}
Large-scale unsupervised language models (LMs) acquire extensive knowledge about the world as well as some reasoning capabilities; however, precisely guiding their outputs remains challenging owing to the wholly unsupervised training paradigm. To enable better control over these models, current approaches typically gather human assessments of the relative quality of outputs and then fine-tune the LM to prefer outputs that align with these judgments, frequently utilizing reinforcement learning from human feedback (RLHF). Yet, RLHF entails a complicated and sometimes unstable workflow: first constructing a reward model that captures human preferences, then applying reinforcement learning to adjust the LM in pursuit of higher estimated reward, all while aiming to preserve the integrity of the original model. This work presents a novel parameterization of the reward model within RLHF that allows for closed-form retrieval of the associated optimal policy. Consequently, we are able to resolve the classical RLHF objective using a straightforward classification loss. The approach, termed \textit{Direct Preference Optimization} (DPO), demonstrates robust performance, stability, and computational efficiency, eliminating both the need for model sampling during fine-tuning and extensive hyperparameter searches. Empirical results indicate that DPO can adjust LMs to human preferences at least as effectively as, and often better than, prior techniques. Particularly, DPO-based fine-tuning outperforms RLHF with PPO in modulating the sentiment of model outputs and matches or surpasses it in quality for tasks like summarization and single-turn dialogue, all while being much easier to implement and train.
\end{abstract}

\section{Introduction}
Large-scale unsupervised language models (LMs) developed on vast corpora have demonstrated unexpected proficiencies~\citep{chowdhery2022palm, brown2020language, touvron2023llama,bubeck2023sparks}. Nevertheless, the datasets used for training are composed of human-generated text reflecting a broad spectrum of objectives, priorities, and expertise. Some of these underlying goals and abilities may be inappropriate for direct emulation. For instance, an AI programming assistant should be able to \textit{recognize} frequent coding errors to effectively correct them; however, during code generation, the model ought to preferentially emulate the (sometimes scarce) examples of expert programming found within its training material. In a parallel vein, although a language model should be capable of identifying a prevalent misconception held by 50\% of individuals, it is undesirable for the model to propagate this false belief in half of its responses regarding that topic. Thus, determining which \emph{responses and behaviors are desirable} from the model’s extensive reservoir of \textit{knowledge and skills} is a fundamental aspect of creating AI that is safe, effective, and responsive to control~\citep{ouyang2022training}. Conventional strategies typically rely on reinforcement learning (RL) to adjust LMs according to human preferences, but as we will demonstrate, the objective optimized in RL-based approaches can in fact be directly attained through a straightforward binary cross-entropy loss, which can considerably streamline the process of preference alignment.

Broadly, current approaches instill targeted behaviors into language models using carefully selected datasets of human preferences, encapsulating the patterns of action deemed secure and beneficial by people. This phase of learning preferences follows an initial, large-scale unsupervised pre-training period using massive textual data. Although the simplest method for aligning with preferences is to perform supervised fine-tuning on curated examples of superior responses from humans, the most advanced approaches belong to the category of reinforcement learning from human or AI feedback (RLHF/RLAIF; \citep{christiano2017deep,bai2022constitutional}). In these frameworks, a reward model is trained on a collection of human preferences, and reinforcement learning is then employed to optimize the LM policy such that it favors outputs with higher assigned rewards, all while constraining the policy from deviating excessively from the pre-trained model. RLHF has enabled the development of models with remarkable prowess in dialog and programming tasks; however, its methodology is notably more involved than supervised techniques. RLHF necessitates training several language models and includes in-training sampling from model policies, substantially increasing computational demands.

This paper introduces a method for optimizing language models to better reflect human preferences, bypassing the need for explicit reward modeling and reinforcement learning techniques. We present \textit{Direct Preference Optimization} (DPO), an algorithm that implicitly targets the same goal as typical RLHF approaches—namely, maximizing rewards subject to a KL-divergence regularization—but offers a more accessible implementation and training process. Conceptually, the DPO mechanism boosts the log likelihood of favored responses compared to less desirable ones, integrating a per-instance adaptive importance weight to counteract the model collapse that can arise from straightforward probability ratio objectives. DPO, like conventional algorithms, uses a formal preference framework (e.g., the Bradley-Terry model; \cite{bradley1952rankanalysis}) for assessing the fit between a reward function and actual human preference data. Unlike standard pipelines that first leverage the preference model to construct and minimize a preference loss for a reward model, and subsequently train a policy to maximize this learned reward, DPO employs a variable substitution to formulate the preference loss directly as a function of the policy itself. As a result, given a set of human preference comparisons over outputs, DPO enables the direct optimization of a policy via a simple binary cross entropy loss, thereby producing an optimal policy with respect to an implicit reward inferred from the observed preferences.

The central contribution of this work is Direct Preference Optimization (DPO), a straightforward, reinforcement learning-free algorithm for training language models based on preference data. Empirical evaluation demonstrates that DPO achieves performance on par with, or better than, established techniques—including PPO-based RLHF—across various tasks, such as sentiment transformation, text summarization, and conversational modeling, in models as large as 6B parameters.

\section{Related Work}

Recent advances in self-supervised language models have shown that increasing model scale allows for zero-shot \citep{radford2019language} and few-shot performance \citep{gpt3,megatron,chowdhery2022palm} on a variety of tasks. Nevertheless, direct fine-tuning on datasets containing explicit instructions and human-authored responses markedly boosts both downstream task performance and alignment with user expectations \citep{mishra-etal-2022-cross,sanh2022multitask,chung2022scaling,thoppilan2022lamda}. This approach, known as `instruction-tuning,' facilitates broader generalization by LLMs to new instructions not represented in the tuning set, thereby enhancing practical applicability \citep{chung2022scaling}. Despite these advances, gathering \textit{relative} human assessments of answer quality—where annotators express preferences rather than provide ideal demonstrations—is often more feasible. For this reason, subsequent research has leveraged datasets of such preferences to further train LLMs, yielding improved results for translation \citep{kreutzer-etal-2018-reliability}, summarization \citep{stiennon2022learning,ziegler2020finetuning}, narrative generation \citep{ziegler2020finetuning}, and adherence to instructions \citep{ouyang2022training,ramamurthy2023is}. In these approaches, a neural reward model is first fit to match the recorded preference data, often using models of comparative judgment such as the Bradley-Terry model \citep{bradley1952rankanalysis}. The language model is then fine-tuned to maximize this reward, employing reinforcement learning algorithms such as REINFORCE \citep{williams1992reinforce}, proximal policy optimization (PPO; \cite{schulman2017proximal}), or their adaptations \citep{ramamurthy2023is}. Parallel research streams use instruction-tuned LLMs, further refined with human feedback, to produce synthetic preference-labeled data targeting properties like safety or non-toxicity \citep{bai2022constitutional}. In these cases, human involvement may be limited to providing rubric-based guidelines for LLM annotation. These developments lie at the intersection of two major areas: reinforcement learning-based training for various objectives in language modeling~\citep{Ranzato2015SequenceLT,paulus2018a,wu2018learning}, and general frameworks for learning from human preferences \citep{christiano2017deep,kupcsik2018learning}. Despite the apparent benefits, applying reinforcement learning to fine-tune large language models with human preferences continues to be a complex and resource-intensive endeavor; in this work, we introduce a principled alternative for optimizing relative preference objectives that circumvents the need for RL.

Beyond linguistic domains, the study of learning policies from preferences has been extensively explored within both the bandit and reinforcement learning frameworks, leading to a variety of proposed methodologies. In contextual bandit scenarios, preferences or rankings are often used in place of explicit rewards—a paradigm referred to as the contextual dueling bandit (CDB; \cite{yue2012karmed,dudik2015contextual}). Without access to absolute reward signals, the theoretical analysis for CDBs replaces the classic notion of optimality with that of a \textit{von Neumann winner}, which is defined as a policy achieving at least a 50\% expected win rate when compared to any alternative policy \citep{dudik2015contextual}. Notably, in CDBs, preference feedback is provided in an online fashion; by contrast, learning from human preferences is generally conducted on a fixed, offline dataset of preference-annotated action pairs \citep{yan2022human}. Likewise, \textit{preference-based RL} (PbRL) considers learning from binary preferences generated by an unspecified 'scoring' function, rather than from explicit reward signals \citep{BusaFekete2014,ruiz2023dueling}. A range of PbRL algorithms have been developed, some of which enable the use of off-policy preference data; however, these typically require an explicit estimation of the latent scoring (reward) function followed by an optimization phase \citep{jain2013learning,BusaFekete2014,christiano2017deep,sadigh2017active,kupcsik2018learning}. In contrast to these two-step pipelines, we propose a unified policy learning strategy that directly adjusts policies to respect observed preferences.

\section{Preliminaries}\label{section:prelims}

We provide an overview of the RLHF workflow as delineated in \citeauthor{ziegler2020finetuning}, and further elaborated in \citep{stiennon2022learning, bai2022training, ouyang2022training}. This process is commonly segmented into three key stages: (1) supervised fine-tuning (SFT); (2) reward modelling via preference collection; and (3) reinforcement learning (RL) based policy optimization.

\textbf{SFT}: The RLHF process typically starts by adapting a pre-trained language model through supervised learning using high-quality datasets relevant to the downstream application (such as dialogue or summarization), thereby yielding an initial policy $\pi^\text{SFT}$.

\textbf{Reward Modelling Phase}: During this subsequent phase, the SFT-tuned model is prompted with inputs $x$, from which pairs of outputs $(y_1, y_2)\sim \pi^\text{SFT}(y \mid x)$ are sampled. Human annotators are then tasked with selecting their preferred completion, resulting in labeled preferences of the form $y_w\succ y_l \mid x$, where $y_w$ and $y_l$ identify the chosen and non-chosen completions, respectively. These preferences are assumed to reflect the judgments of an underlying, unobserved reward model $r^*(y, x)$. Several models are available to capture this preference structure; among them, the Bradley-Terry (BT) model \cite{bradley1952rankanalysis} is commonly used (though broader frameworks such as the Plackett-Luce model \citep{plackett1975analysis, luce2012individual} can also be employed if multiple answer rankings are accessible). In the BT framework, the distribution over human preferences, $p^*$, is formalized as follows:

\begin{equation}\label{eq:bradley-terry}
    p^*(y_1\succ y_2 \mid x)=\frac{\exp\left(r^*(x, y_1)\right)}{\exp\left(r^*(x, y_1)\right) + \exp\left(r^*(x, y_2)\right)}.
\end{equation}

Given a fixed dataset of comparative examples $\mathcal{D}=\bigl\{x^{(i)}, y_w^{(i)}, y_l^{(i)}\bigr\}_{i=1}^N$ sampled according to $p^*$, one can model a reward function $r_{\phi}(x, y)$ and estimate the parameter $\phi$ via maximum likelihood estimation. By recasting this task as a binary classification problem, the associated negative log-likelihood loss can be written as:

\begin{equation}\label{eq:reward_model}
    \mathcal{L}_R(r_{\phi}, \mathcal{D}) = -\mathbb{E}_{(x, y_w, y_l)\sim \mathcal{D}}\bigl[\log \sigma(r_{\phi}(x, y_w)- r_{\phi}(x, y_l))\bigr]
\end{equation}

where $\sigma$ denotes the logistic sigmoid function. For language models, $r_{\phi}(x, y)$ typically inherits weights from the SFT model $\pi^\text{SFT}(y \mid x)$, with an added linear transformation after the last transformer block, producing a scalar score as reward \cite{ziegler2020finetuning}. To reduce variance in the reward function, it is common to apply normalization such that $\mathbb{E}_{x,y\sim \mathcal{D}}\left[r_\phi(x, y)\right] = 0$ for all $x$.

\textbf{RL Fine-Tuning Phase}: Within the RL fine-tuning stage, the learned reward model furnishes supervision to guide the language model’s policy. Following established approaches~\citep{jaques2017sequence, jaques2020human}, the optimization objective takes the form:

\begin{equation}\label{eq:RL}
\max_{\pi_{\theta}}  \mathbb{E}_{x\sim \mathcal{D}, y\sim \pi_{\theta}(y \mid x)}\bigl[r_{\phi}(x, y)\bigr] - \beta\mathbb{D}_{\textrm{KL}}\bigl[\pi_{\theta}(y\mid x)\mid \mid \pi_\text{ref}(y\mid x)\bigr],
\end{equation}

where $\beta$ controls regularization to stay close to a reference policy $\pi_\text{ref}$, which is typically the SFT model $\pi^\text{SFT}$. Usually, the policy $\pi_\theta$ is initialized to match $\pi^\text{SFT}$. This regularization is essential to keep the model’s outputs within the reward model's support, sustain generation diversity, and avoid collapsing to singular high-reward outputs. As language output is discrete and the objective is not directly differentiable, it is generally tackled with reinforcement learning techniques. The prevalent strategy \citep{ziegler2020finetuning, stiennon2022learning, bai2022training, ouyang2022training} is to define the reward as ${r(x, y) = r_{\phi}(x, y) -\beta (\log \pi_{\theta}(y\mid x) - \log \pi_\text{ref}(y\mid x))}$ and maximize this surrogate reward using PPO \cite{schulman2017proximal}.

\section{Direct Preference Optimization}\label{sec:DPO}

Seeking to address the complexities of reinforcement learning in large-scale applications such as language model fine-tuning, we aim to devise a more straightforward preference-driven policy optimization technique. In contrast to standard RLHF frameworks that construct an explicit reward to be optimized through RL, we consider a reward model parameterization where the optimal policy can be analytically determined—thereby eliminating the need for a reinforcement learning loop. As we elaborate below, the central idea is to utilize an exact mapping from reward to optimal policy, which allows reformulating the loss on rewards into a loss directly over policies. This change-of-variables procedure bypasses the need to fit a distinct reward model while ensuring that optimization remains aligned with common models of human preference such as the Bradley-Terry model. Ultimately, the policy network serves the dual purpose of representing both the generative model and the underlying (implicit) reward signal.

\textbf{Deriving the DPO objective.} We begin by considering the same reinforcement learning objective stated in Eq.~\ref{eq:RL}, which operates under a general reward function $r$. As established in earlier studies~\citep{peters2007reinforcement, peng2019advantage, korbak2022reinforcement, go2023aligning}, it can be readily demonstrated that the solution optimizing the KL-regularized reward maximization problem described in Eq.~\ref{eq:RL} adopts the following form:

\begin{equation}\label{eq:op_policy}
    \pi_r(y\mid x) = \frac{1}{Z(x)}\pi_\text{ref}(y\mid x)\exp\left(\frac{1}{\beta}r(x, y)\right),
\end{equation}

where $Z(x) =\sum_{y}\pi_\text{ref}(y\mid x)\exp\left(\frac{1}{\beta}r(x, y)\right)$ denotes the partition function. A detailed derivation is available in Appendix \ref{app:derivation1}. Even with an MLE-based estimate $r_{\phi}$ for the underlying reward function $r^*$, computing the partition function $Z(x)$ remains computationally intensive \citep{korbak2022reinforcement, go2023aligning}, rendering this policy representation challenging for practical use. Nevertheless, it is possible to reformulate Eq.~\ref{eq:op_policy} to represent the reward in terms of the corresponding optimal policy $\pi_r$, the reference policy $\pi_\text{ref}$, and the unknown $Z(\cdot)$. By taking the logarithm of both sides of Eq.~\ref{eq:op_policy} and rearranging terms, we derive:

\begin{equation}\label{eq:main_eq}
    r(x,y) =\beta \log \frac{\pi_r(y\mid x)}{\pi_\text{ref}(y\mid x)} + \beta \log Z(x).
\end{equation}

This reparameterization is also valid for the true reward $r^*$ and its associated optimal policy $\pi^*$. Importantly, the Bradley-Terry model’s structure relies solely on the difference between the rewards assigned to two outputs; namely, ${p^*(y_1 \succ y_2 \mid x) = \sigma(r^*(x, y_1) - r^*(x, y_2))}$. Plugging the result from Eq.~\ref{eq:main_eq} for $r^*(x,y)$ into the preference model from Eq.~\ref{eq:bradley-terry} leads to a cancellation of the partition function, enabling an expression of the probability of human preference strictly through $\pi^*$ and $\pi_\text{ref}$. Thus, under the Bradley-Terry model, the optimal RLHF policy $\pi^*$ adheres to the following preference formulation:

\begin{equation}\label{eq:objective}
    p^*(y_1\succ y_2 \mid x)=\frac{1}{1 + \exp\left(\beta \log \frac{\pi^*(y_2\mid x)}{\pi_\text{ref}(y_2\mid x)} - \beta \log \frac{\pi^*(y_1\mid x)}{\pi_\text{ref}(y_1\mid x)}\right)}
\end{equation}

Further mathematical details are included in Appendix~\ref{app:derivation2}. While Eq.~\ref{eq:objective} specifically addresses the Bradley-Terry model, parallel results can be obtained for broader Plackett-Luce models~\citep{plackett1975analysis, luce2012individual}, as demonstrated in Appendix~\ref{app:plackett_luce_models}.

With the human preference likelihood now characterized by the optimal policy rather than the reward, we may now design a maximum likelihood estimation objective for a parameterized policy $\pi_\theta$. Similar to reward modeling (see Eq.~\ref{eq:reward_model}), our policy objective is expressed as:

\begin{equation}\label{eq:optimum_model}
    \mathcal{L}_\text{DPO}(\pi_{\theta}; \pi_\text{ref}) = -\mathbb{E}_{(x, y_w, y_l)\sim \mathcal{D}}\left[\log \sigma \left(\beta \log \frac{\pi_{\theta}(y_w\mid x)}{\pi_\text{ref}(y_w\mid x)} - \beta \log \frac{\pi_{\theta}(y_l\mid x

In this framework, we learn an implicit reward through an alternative formulation, for which the optimal policy coincides with $\pi_\theta$. Notably, our approach is tantamount to fitting a reparameterized Bradley-Terry model, which imparts certain desirable theoretical guarantees, such as consistency under appropriate assumptions regarding the distribution of preference data \cite{bong2022generalized}. We elaborate on the theoretical aspects of DPO, including connections to related literature, in Section~\ref{sec:theory}.

\textbf{How does the DPO update operate?} To gain mechanistic insight into DPO, it is helpful to inspect the gradient of its loss function $\mathcal{L}_\text{DPO}$. The gradient with respect to model parameters $\theta$ takes the following form:

\begin{multline*}\label{eq:gradient}
    \nabla_\theta \mathcal{L}_\text{DPO}(\pi_\theta;\pi_\text{ref}) = \\ -\beta\mathbb{E}_{(x, y_w, y_l) \sim \mathcal{D}} \bigg[\underbrace{\sigma(\hat{r}_\theta(x, y_l) - \hat{r}_\theta (x, y_w))}_\text{increased influence when reward estimate is inaccurate}\bigg[\underbrace{\nabla_\theta\log \pi(y_w \mid x)}_\text{push up $y_w$ likelihood} - \underbrace{\nabla_\theta\log\pi(y_l \mid x)}_\text{push down $y_l$ likelihood}\bigg]\bigg],
\end{multline*}

Here, $\hat{r}_\theta(x, y) = \beta \log \frac{\pi_\theta(y \mid x)}{\pi_\text{ref}(y \mid x)}$ represents the reward inferred implicitly by $\pi_\theta$ relative to the reference policy $\pi_\text{ref}$ (see further discussion in Section~\ref{sec:theory}). Conceptually, the gradient encourages higher probability for the preferred outputs $y_w$ and suppresses the probability of less-preferred outputs $y_l$. Crucially, each training instance is weighted according to the degree to which the implicit reward $\hat{r}_\theta$ erroneously assigns greater value to the dispreferred $y_l$, with scaling determined by $\beta$. This reflects the magnitude of misordering by the implicit reward, tempered by the KL-divergence constraint strength. Empirical results highlight the necessity of this weighting; omitting the weighting coefficient (as in a na\"ive implementation) can lead the language model to degrade in performance (see Appendix Table~\ref{tab:unlikelihood_generations}).

\textbf{DPO overview.} The standard DPO workflow consists of two primary stages: (1) For each input prompt $x$, generate response samples $y_1, y_2 \sim \pi_\text{ref}(\cdot \mid x)$, and use human preference annotations to build an offline preference dataset $\mathcal{D} = \{x^{(i)}, y_w^{(i)}, y_l^{(i)}\}_{i=1}^N$; (2) Train the language model $\pi_\theta$ by minimizing $\mathcal{L}_\text{DPO}$, using the designated $\pi_\text{ref}$, dataset $\mathcal{D}$, and temperature parameter $\beta$. In practice, to avoid the need for collecting new data, publicly released preference datasets are typically employed rather than sampling new outputs and eliciting human choices. Since these preference datasets are generated from $\pi^\text{SFT}$, it is preferable to set $\pi_\text{ref} = \pi^\text{SFT}$ if such a model exists. In the absence of $\pi^\text{SFT}$, $\pi_\text{ref}$ is instead constructed by maximizing the likelihood of the preferred completions ${(x, y_w)}$, namely, ${\pi_\text{ref} = \argmax_{\pi}\mathbb{E}_{x, y_w \sim \mathcal{D}}\left[\log \pi(y_w \mid x)\right]}$. This strategy serves to bridge the gap caused by distribution shift between the real but inaccessible reference distribution and the practical $\pi_\text{ref}$ deployed in DPO. Implementation specifics and details about hyperparameter settings are further discussed in Appendix~\ref{app:implementation}.

\section{Theoretical Analysis of DPO} This section develops further intuition and theoretical foundations for DPO, elucidates its properties, and explains how its advantages address limitations inherent to actor-critic methods commonly used for RLHF (for example, PPO~\cite{schulman2017proximal}).

\label{sec:theory}

\subsection{Your Language Model Is Secretly a Reward Model} DPO is distinctive in that it unifies learning a policy and optimizing for preferences into a single maximum likelihood framework, eliminating the need to separately fit an explicit reward function or employ RL to train the policy. Importantly, the optimization objective in Eq. \ref{eq:main_eq} mirrors the Bradley-Terry model, with the reward expressed as $r^*(x, y) = \beta \log\frac{\pi^*_\theta(y \mid x)}{\pi_\text{ref}(y \mid x)}$. The parametric model $\pi_{\theta}$ is thus updated analogously to the reward model in Eq. \ref{eq:reward_model}, given a reparameterization. In what follows, we present the formal theory justifying this approach, demonstrate that it does not reduce the expressive power of possible reward models, and prove that this construction enables exact retrieval of the optimal policy. We start by introducing an equivalence relation for reward functions.

\begin{definition}
We define two reward functions $r(x, y)$ and $r'(x, y)$ to be equivalent if and only if there exists a function $f$ such that ${r(x, y)-r'(x, y) = f(x)}$.     
\end{definition}

One can readily verify that this equivalence definition indeed constitutes an equivalence relation, thereby dividing the collection of reward functions into equivalence classes. This leads to the following two lemmas:

\begin{lemma}\label{lemma:same_prefrence} Within the Plackett-Luce framework, including the Bradley-Terry model, any pair of reward functions that belong to the same equivalence class generates identical preference distributions.
\end{lemma}

\begin{lemma}\label{lemma:same_policy}
    If two reward functions are members of the same equivalence class, they induce the same optimal policy in the context of the constrained RL formulation.
\end{lemma}

As the arguments are straightforward, we relegate their proofs to Appendix \ref{app:lemma1}. The first lemma highlights a classic issue of under-specification characteristic of the Plackett-Luce model family \cite{plackett1975analysis}. This ambiguity typically necessitates the introduction of identifiability restrictions in order to ensure well-posed MLE estimation from Eq. \ref{eq:reward_model} \cite{bong2022generalized}. The second lemma asserts that all reward functions within the same equivalence class result in identical optimal policies; thus, for the purposes of optimization, it suffices to recover any representative of the corresponding equivalence class. The following theorem, with its proof in Appendix~\ref{app:thm1}, formalizes this notion:

\begin{theorem}\label{thm:main}
    Subject to mild assumptions, every reward class that is consistent with the Plackett-Luce (including Bradley-Terry) models can be realized through the parameterization ${r(x, y) = \beta \log \frac{\pi(y\mid x)}{\pi_\text{ref}(y\mid x)}}$ for some model $\pi(y\mid x)$ and a specified reference model $\pi_\text{ref}(y \mid x)$.
\end{theorem}

\begin{sproof}
    Take any given reward function $r(x, y)$, which corresponds to some optimal model $\pi_r(y \mid x)$ defined by Eq. \ref{eq:op_policy}. We aim to show that a member of the equivalence class of $r$ can be written in the proposed parameterized form. Specifically, we construct the projection operator $f$ as  
\begin{equation}
    f(r; \pi_\text{ref}, \beta)(x, y) = r(x, y) - \beta\log\sum_{y}\pi_\text{ref}(y\mid x)\exp\left(\frac{1}{\beta}r(x, y)\right)
\end{equation}
This operator adjusts the reward function by subtracting the log-partition function, which depends solely on the context $x$. As this normalization term is a function of $x$ alone, the transformed reward $f(r; \pi_\text{ref}, \beta)(x, y)$ remains in the same equivalence class as $r(x, y)$. Substituting $r$ with the right-hand side of Eq.~\ref{eq:main_eq}, which is valid for any reward, yields $f(r; \pi_\text{ref}, \beta)(x, y) = \beta \log \frac{\pi_r(y\mid x)}{\pi_\text{ref}(y\mid x)}$. Thus, $f$ generates an element in the equivalence class of $r$ with the required form, confirming that the chosen reparameterization is fully representative for our purposes.
\end{sproof}

An alternative interpretation of Theorem~\ref{thm:main} is that it delineates precisely which reward function, among each equivalence class, is chosen by the DPO reparameterization—namely, the reward function that satisfies

\begin{equation}\label{eq:lag_p}
     \sum_{y}\underbrace{\pi_\text{ref}(y\mid x)\exp\left(\frac{1}{\beta}r(x, y)\right)}_{=\pi(y\mid x)\text{, using Thm.~\ref{thm:main} reparam.}} = 1,
\end{equation}

which ensures that $\pi(y\mid x)$ forms a valid probability distribution, meaning that its values are non-negative and sum to unity. Furthermore, as indicated by Eq.~\ref{eq:op_policy}, Eq.~\ref{eq:lag_p} represents the partition function of the optimal policy derived from the reward function $r(x, y)$. The central contribution of the DPO algorithm lies in the imposition of specific constraints on the otherwise unconstrained Plackett-Luce family of preference models (and, in particular, the Bradley-Terry model). These constraints retain the expressive capacity of the reward model class, while ensuring that the optimal policy given by Eq.~\ref{eq:op_policy} admits a closed-form solution across all inputs $x$.

\subsection{Instability of Actor-Critic Algorithms} Our framework further enables a diagnosis of the instability issues that can arise with standard actor-critic methods commonly adopted in RLHF, such as PPO. We proceed in accordance with the RLHF workflow and center our attention on the RL fine-tuning process described in Section \ref{section:prelims}. This perspective reveals links to the control as inference approach \cite{levine2018reinforcement}, specifically for the constrained RL formulation described in \ref{eq:RL}. Assuming a parameterized policy $\pi_{\theta}(y\mid x)$, the objective becomes the minimization of $\mathbb{D}_{\text{KL}}[\pi_{\theta}(y|x) \mid \mid \pi^*(y\mid x)]$, where $\pi^*$, given by Eq. \ref{eq:optimum_model}, denotes the optimal policy defined in terms of the reward function $r_{\phi}(y, x)$. After some mathematical manipulation, the resulting optimization target is:

\begin{equation}\label{eq:AC}
    \max_{\pi_{\theta}}\mathbb{E}_{\pi_{\theta}(y\mid x)}\bigg[\underbrace{r_{\phi}(x, y) -\beta\log\sum_{y}\pi_\text{ref}(y\mid x)\exp\left(\frac{1}{\beta}r_{\phi}(x, y)\right)}_{f(r_{\phi}, \pi_\text{ref}, \beta)} - \underbrace{\beta\log\frac{\pi_{\theta}(y\mid x)}{\pi_\text{ref}(y\mid x)}}_{\text{KL}}\bigg]
\end{equation}

The objective here mirrors that employed in previous research 
\citep{ziegler2020finetuning, stiennon2022learning, bai2022training, ouyang2022training}, leveraging the DPO-equivalent reward formulated for the reward class $r_{\phi}$. Under this formulation, the normalization factor within $f(r_{\phi}, \pi_\text{ref}, \beta)$ can be seen as representing the soft value function corresponding to the reference policy $\pi_\text{ref}$. Although this normalization does not influence the location of the optimum, omitting it may result in policy gradients with substantial variance, thereby impeding stable optimization. It is possible to account for this normalization term through the introduction of a learned value function, but optimizing such a function introduces its own set of challenges. Another approach found in the literature involves normalizing the rewards relative to a human completion baseline—effectively approximating the normalization with a Monte Carlo estimate based on a single sample. By contrast, the DPO reparameterization provides a reward formulation that is baseline-free.

\section{Experiments}
Here, we empirically assess how effectively DPO can train policies based exclusively on preference data. We begin with controlled experiments in text generation to compare how efficiently DPO manages the trade-off between maximizing the reward and constraining the KL-divergence from the reference policy, relative to established methods such as PPO. Subsequently, we extend our evaluation to encompass larger models and more challenging RLHF tasks, including those involving summarization and dialog. Our findings suggest that, with minimal hyperparameter tuning, DPO matches or surpasses competitive baselines, such as RLHF using PPO and strategies that select the best of $N$ trajectories as determined by a learned reward model. Prior to presenting these empirical results, we outline the experimental procedures; further specifics can be found in Appendix~\ref{app:exp_details}.

\textbf{Tasks.} Our study examines three distinct tasks within the domain of open-ended text generation. Across all tasks, the goal is to learn a policy from a preference dataset $\mathcal{D} = \bigl\{x^{(i)}, y_w^{(i)}, y_l^{(i)}\bigr\}_{i=1}^N$. For the \textbf{controlled sentiment generation} task, $x$ consists of an initial segment of a movie review drawn from the IMDb dataset \cite{maas-EtAl:2011:ACL-HLT2011}, with the policy required to continue the text $y$ so that it conveys a positive sentiment. To facilitate rigorous assessment, we automatically construct preference pairs between generated completions using a sentiment classifier, ensuring that $p(\text{positive} \mid x, y_w) > p(\text{positive} \mid x, y_l)$. The SFT approach here involves further training GPT-2-large to convergence using the training portion of IMDb reviews (see Appendix~\ref{app:sentiment_details} for more specifics).

For the \textbf{summarization} task, $x$ is a Reddit forum post, and the objective is to generate a concise summary $y$ highlighting its principal arguments. We make use of the Reddit TL;DR dataset \citep{volske-etal-2017-tl} and incorporate human-labeled preferences collected by \citeauthor{stiennon2022learning}. Our SFT baseline is produced by fine-tuning a language model on a set of summaries composed by humans,\footnote{\url{https://huggingface.co/CarperAI/openai_summarize_tldr_sft}} employing the TRLX \citep{leandro_von_werra_2023_7790115} toolkit for RLHF training. The human preferences correspond to outputs from a related, independently fine-tuned SFT model as curated by \citeauthor{stiennon2022learning}.

Lastly, in the context of \textbf{single-turn dialogue}, $x$ is a user-issued prompt, which can range from scientific queries to personal advice-seeking. The policy is expected to deliver a relevant and useful reply $y$. For this task, we utilize the Anthropic Helpful and Harmless dialogue corpus \citep{bai2022training}, which comprises 170,000 exchanges involving human participants and an automated assistant. Each interaction concludes with two responses from a high-capacity (but unspecified) language model, accompanied by a label indicating which response the human judge preferred. Since a pre-trained SFT model is unavailable for this dataset, we instead adapt an existing language model through fine-tuning solely on the preferred completions to create the SFT reference.

\textbf{Evaluation.} We employ two distinct strategies for assessing our algorithms. For the constrained reward maximization task in controlled sentiment generation, we evaluate each approach by mapping the trade-off frontier between realized reward and KL-divergence from the reference policy, which is feasible due to access to an explicit ground-truth reward function (a sentiment classifier). In contrast, since real-world applications lack access to ground-truth rewards, we instead report the \textit{win rate} of each algorithm relative to a baseline policy. For this, we use GPT-4 as an automated evaluator serving as a surrogate for human judgment: in the summarization context, summary quality is compared, while for single-turn dialogue, the focus is on helpfulness. The baseline in summarization consists of gold-standard reference summaries from the test set, and in dialogue, we use the preferred responses annotated within the test data. While there is evidence that language models can outperform traditional metrics as automatic evaluators \citep{Chen2023ExploringTU}, we bolster our reliance on GPT-4 by performing a supplementary human evaluation, as detailed in Sec.~\ref{sec:human-judgments}. Results reveal strong concordance between GPT-4 and human judgments, with rates of agreement between GPT-4 and human annotators comparable to, or exceeding, inter-human reliability.

\textbf{Methods.} Beyond DPO, we benchmark a variety of established techniques for aligning language models to human preference signals. At the simplest level, we consider zero-shot prompting using \textbf{GPT-J} \citep{gpt-j} for summarization, and 2-shot prompting with \textbf{Pythia-2.8B} \citep{biderman2023pythia} for dialogue. Further, we test the \textbf{SFT} baseline, as well as \textbf{Preferred-FT}, which involves supervised fine-tuning on preferred completions $y_w$—sourced either from the SFT model (in the sentiment and summarization tasks) or from a generic LM (in dialogue). An alternative supervised approach, \textbf{Unlikelihood}~\citep{welleck2019neural}, seeks to maximize the predicted probability of $y_w$ while simultaneously minimizing that of $y_l$, introducing a tunable `unlikelihood' penalty via an optional coefficient $\alpha\in[0,1]$. Additionally, we include \textbf{PPO} \citep{schulman2017proximal}, trained with a reward model based on preference data, as well as \textbf{PPO-GT}, which is an idealized variant leveraging the known ground-truth reward function within the controlled sentiment scenario. For sentiment-based experiments, PPO-GT is instantiated in two ways: using an out-of-the-box implementation \cite{leandro_von_werra_2023_7790115}, and a custom version that normalizes reward signals and tunes hyperparameters for enhanced performance; the same enhancements are also applied to PPO trained with learned rewards. We further assess a \textbf{Best of $N$} strategy: here, $N$ outputs are sampled from the SFT model (or from Preferred-FT in dialogue), and the response receiving the highest reward (per the preference-trained model) is returned. Although this method is often strong—decoupling the reward model’s efficacy from PPO optimization—its computational demands are prohibitive for even moderate $N$, as each test-time input requires sampling $N$ separate completions.

\subsection{How well can DPO optimize the RLHF objective?}

The KL-constrained reward maximization objective commonly employed in RLHF algorithms seeks to strike a balance between maximizing reward and maintaining proximity to the reference policy. Thus, evaluating and comparing such algorithms necessitates considering both the obtained reward and the KL divergence from the reference; attaining marginally higher rewards at the expense of substantially greater KL divergence may not be a favorable outcome. Figure~\ref{fig:frontier-tldr-main} illustrates the tradeoff between reward and KL divergence for different methods within the sentiment analysis task. For each algorithm, we conduct several training runs, each using a unique value for the policy conservativeness hyperparameter (specifically, target KL $\in\{3,6,9,12\}$ for PPO, $\beta \in \{0.05,0.1,1,5\}$, $\alpha\in\{0.05,0.1,0.5,1\}$ for unlikelihood, and different random seeds for preferred-FT), amounting to a total of 22 distinct training configurations. At intervals of every 100 training steps until the models reach convergence, we assess each policy across a standardized test set, measuring the mean reward as determined by the actual reward function and computing the average sequence-level KL\footnote{That is, the sum of the per-timestep KL-divergences.} relative to the reference policy $\text{KL}\left(\pi\mid \mid \pi_\text{ref}\right)$. Our findings indicate that DPO establishes the most favorable reward-KL tradeoff frontier, securing the highest rewards without incurring large KL deviations. This observation is striking for several reasons. Primarily, although DPO and PPO are designed to optimize the same objective, DPO accomplishes this with superior efficiency—its performance in the reward versus KL space consistently exceeds that of PPO. Additionally, DPO surpasses PPO in constructing this efficient frontier, \emph{even in scenarios where PPO has direct access to the true reward signal} (PPO-GT).

\subsection{Can DPO scale to real preference datasets?}
\label{sec:dpo-real-datasets}
We proceed to assess the effectiveness of DPO when fine-tuned on real-world preference datasets for both summarization and single-turn dialogue tasks. In the context of summarization, standard automated metrics like ROUGE often do not align closely with human judgment~\citep{stiennon2022learning}, and it has been previously demonstrated that models fine-tuned with PPO on human preference data can generate superior summaries. To compare various approaches, we sample model outputs from the test split of the TL;DR summarization dataset and measure their mean win rate versus reference outputs in the test set. For all evaluated methods, completions are drawn at temperature settings ranging from 0.0 to 1.0, and the corresponding win rates are presented in Figure~\ref{fig:frontier-tldr-main} (right). DPO, PPO, and Preferred-FT are each initialized from the same GPT-J SFT model\footnote{\url{https://huggingface.co/CarperAI/openai_summarize_tldr_sft}}. Our findings indicate that DPO attains an approximate win rate of 61\% at a temperature of 0.0, surpassing PPO, which achieves a win rate of roughly 57\% at its own optimal temperature of 0.0. Additionally, DPO secures a higher best-case win rate than the best of $N$ baseline. Notably, we did not undertake extensive tuning of the $\beta$ hyperparameter for DPO, suggesting that these outcomes may not reflect the full capabilities of the method. Furthermore, DPO exhibits significantly greater resilience to changes in sampling temperature compared to PPO, whose performance declines to that of the unmodified GPT-J model as temperature increases. Preferred-FT offers little to no improvement over the SFT baseline. In Section~\ref{sec:human-judgments}, we also conduct a direct comparison between DPO and PPO in human evaluation, finding that DPO samples generated at temperature 0.25 are preferred 58\% of the time over those from PPO at temperature 0.

For single-turn dialogue evaluation, we assess the various approaches using the portion of the Anthropic HH dataset \citep{bai2022training} test split that contains a single exchange between human and assistant. In our GPT-4 based assessments, the methods’ win rates are measured with respect to the preferred test completions as the ground truth. Since there is no established SFT baseline for this benchmark, our process begins with the pre-trained Pythia-2.8B model; we apply Preferred-FT to adapt a reference model on the selected completions, ensuring the outputs remain in-distribution, followed by training with DPO. Additionally, we benchmark against the top completion out of 128 from Preferred-FT runs (empirically, the performance of the Best of $N$ approach saturates at $N=128$ for this task; see Appendix Figure~\ref{fig:best-of-n}), as well as a 2-shot prompt version of the base Pythia-2.8B. Across the methods and optimal temperature choices, DPO matches or exceeds the best alternatives. We also analyze an RLHF model fine-tuned via PPO on Anthropic HH \footnote{\url{https://huggingface.co/reciprocate/ppo_hh_pythia-6B}} with a reputable implementation \footnote{\url{https://github.com/CarperAI/trlx/tree/main/examples/hh}}, but are unable to identify any prompting or sampling configuration that yields superior results to the original Pythia-2.8B. Drawing from findings in TL;DR and recognizing that both the PPO and Best of 128 approaches aim to maximize the same reward, we use Best of 128 as an approximate stand-in for PPO performance. Notably, DPO is unique among efficient methods in consistently surpassing the preferred completions from the Anthropic HH dataset, achieving comparable or even better outcomes than the more compute-intensive Best of 128 baseline. Figure~\ref{fig:dialogue-main} further demonstrates that DPO attains peak effectiveness in relatively few training steps.

\subsection{Generalization to a new input distribution}

To further analyze how PPO and DPO handle distributional changes, we assess the policies trained in our Reddit TL;DR summarization setup by applying them to a distinct domain: the test split of the CNN/DailyMail dataset \citep{nallapati-etal-2016-abstractive}, which comprises news articles. We employ the optimal sampling temperatures from the TL;DR task (0 and 0.25) for both algorithms. The evaluation metrics are reported in Table~\ref{tab:ood}. The GPT-4 win rate is computed by comparing model outputs to reference summaries in these datasets, using a prompt adapted from the one used in the Reddit TL;DR experiment (substituting “forum post” with “news article”). On this shifted distribution, DPO maintains a considerable lead over PPO. These findings offer preliminary support that DPO-trained policies can generalize comparably well to those trained via PPO, despite DPO not leveraging the extra unlabeled Reddit TL;DR prompts available to PPO.

\subsection{Validating GPT-4 judgments with human judgments}
\label{sec:human-judgments}

To assess the trustworthiness of GPT-4’s preference judgments, we conduct a human evaluation on outputs from the TL;DR summarization task, using two variants of the GPT-4 prompt. The first, \textbf{GPT-4 (S)} (simple), asks GPT-4 to select the summary best capturing the main points of the post. The second, \textbf{GPT-4 (C)} (concise), also directs GPT-4 to consider conciseness, motivated by observations that the simple prompt tends to make GPT-4 favor summaries that are lengthier and more redundant compared to human raters. Full prompt text is available in Appendix~\ref{app:prompts}. We compare three system outputs with varying quality: the highest-scoring (DPO at temperature 0.25), the lowest-scoring (PPO at temperature 1.0), and an intermediate-quality model (SFT at temperature 0.25). Each is judged relative to PPO sampled greedily (its best-performing configuration), aiming to span a broad quality spectrum. In both prompt settings, GPT-4’s choices align with human rater choices at approximately the same frequency that humans concur with one another, implying that GPT-4 is a reasonable stand-in for direct human assessment (note that for DPO and PPO-1, multiple human annotations are collected, but for others, human annotator resources were limited). Overall, the \textbf{GPT-4 (C)} prompt produces win rates more closely matching human judgments, so it is adopted for primary analysis in Section~\ref{sec:dpo-real-datasets}. Further methodological details, including the rating interface and human volunteer information, are presented in Appendix~\ref{app:human-study}.

\section{Discussion}
Preference-based learning offers an effective and scalable approach to developing capable and aligned language models. In this work, we have presented DPO, a straightforward method for training language models from preference data that does not rely on reinforcement learning. Instead of adapting the preference learning problem into a conventional RL framework to leverage existing RL algorithms, DPO establishes a direct correspondence between language model policies and reward functions. This mapping makes it possible to train models to adhere to human preferences using only a basic cross-entropy loss function, bypassing both reinforcement learning procedures and any compromise in model performance. Remarkably, with minimal hyperparameter tuning, DPO achieves results that are on par with or superior to traditional RLHF methods such as those utilizing PPO, thereby lowering the entry barrier to preference-driven language model training.

\textbf{Limitations \& Future Work.} Our findings give rise to multiple avenues for further exploration. One open question concerns the generalization capabilities of DPO-trained policies on out-of-distribution data in comparison to models trained with explicit reward functions. Preliminary evidence indicates that DPO can match the generalization ability of PPO-based approaches, but a more extensive evaluation is warranted. Furthermore, it is worth investigating whether employing DPO’s policy for self-labeling can be leveraged effectively for unlabeled prompt datasets. Another interesting question is how over-optimization of the reward might present itself within direct preference optimization, and whether the minor dip in performance observed in Figure~\ref{fig:dialogue-main}-right could be attributed to this phenomenon. Moreover, though our analysis encompasses models up to 6B parameters, scaling DPO to substantially larger, state-of-the-art language models represents a promising area for future research. From an evaluation standpoint, we have observed that the prompt content influences GPT-4-derived win rates, highlighting the need for studies focused on optimal methods for obtaining reliable assessments from automated evaluators. Lastly, there is substantial potential for applying DPO beyond human preference-based language model training, such as in the optimization of generative models for different data modalities.

\end{document}